\section{Conclusiones}

El análisis desarrollado permitió comprender la diversidad de enfoques existentes en el diseño e implementación de sistemas operativos, evidenciando que no existe una única arquitectura óptima, sino soluciones adaptadas a distintos propósitos, contextos y niveles de complejidad. La comparación entre sistemas monolíticos, microkernel, híbridos y experimentales facilitó la identificación de las ventajas y limitaciones de cada propuesta, así como su impacto en aspectos clave como el rendimiento, la seguridad, la estabilidad y la mantenibilidad.

Asimismo, el estudio de sistemas operativos consolidados y alternativos puso de manifiesto la importancia de la documentación, la comunidad de desarrollo y las herramientas disponibles como factores determinantes para su adopción en entornos académicos. Los sistemas con documentación clara, código bien estructurado y soporte activo ofrecen mayores oportunidades de aprendizaje, permitiendo a los estudiantes analizar internamente componentes críticos como la gestión de procesos, memoria, sistemas de archivos y mecanismos de comunicación.

La evaluación comparativa realizada permitió justificar la selección de SerenityOS como sistema de referencia para el desarrollo del trabajo, debido a su equilibrio entre complejidad técnica y valor educativo. Su arquitectura bien definida, el uso predominante de C++ moderno, la modularidad del código y la disponibilidad de herramientas de compilación, ejecución y depuración lo convierten en una plataforma adecuada para el estudio práctico de los conceptos fundamentales de sistemas operativos.

Finalmente, este trabajo establece una base sólida para futuras etapas de análisis y desarrollo, orientadas a la experimentación, extensión o implementación de nuevos módulos sobre el sistema seleccionado. La metodología aplicada y los criterios definidos pueden reutilizarse en investigaciones posteriores, contribuyendo a una comprensión más profunda del funcionamiento interno de los sistemas operativos y fortaleciendo la formación académica en esta área fundamental de la ingeniería informática.
